\documentclass{article}
\usepackage{amsmath,amssymb,amsthm,algorithm,algorithmic}
\usepackage{geometry}
\geometry{margin=1in}
\newtheorem{theorem}{Theorem}
\newtheorem{lemma}{Lemma}
\newcommand{\prox}{\operatorname{prox}}
\newcommand{\inner}[2]{\langle #1,#2\rangle}
\newcommand{\norm}[1]{\|#1\|}
\begin{document}

\section*{Primal–Dual Hybrid Gradient (PDHG) Algorithm}

\section*{Mathematical Formulation}
Consider the convex–concave saddle‑point problem
\[
\min_{x\in\mathbb R^{n}} \; f(x)+g(Kx),
\qquad 
\text{equivalently}\qquad 
\min_{x}\max_{y}\; \mathcal L(x,y):=f(x)+\inner{Kx}{y}-g^{*}(y),
\]
where 
\begin{itemize}
  \item $f:\mathbb R^{n}\to\mathbb R$ is convex, differentiable and $L$‑smooth,
  \item $g:\mathbb R^{m}\to\mathbb R\cup\{+\infty\}$ is convex (possibly nonsmooth),
  \item $K\in\mathbb R^{m\times n}$ is a linear operator,
  \item $g^{*}$ is the convex conjugate of $g$.
\end{itemize}
The PDHG iterates are
\begin{align}
y_{k+1}&=\prox_{\sigma g^{*}}\!\bigl(y_{k}+\sigma K\bar x_{k}\bigr), \label{eq:dual}\\[4pt]
x_{k+1}&=\prox_{\tau f}\!\bigl(x_{k}-\tau K^{\top}y_{k+1}\bigr), \label{eq:primal}\\[4pt]
\bar x_{k+1}&=x_{k+1}+\theta\,(x_{k+1}-x_{k}), \qquad \theta\in[0,1]. \label{eq:extrap}
\end{align}
Parameters $\tau>0$, $\sigma>0$ are step sizes and must satisfy the
\[
\tau\sigma\|K\|^{2}<1, \qquad\text{(stability condition)} \tag{C1}
\]
where $\|K\|$ denotes the spectral norm of $K$.

When $g\equiv0$, $\prox_{\sigma g^{*}}$ reduces to the identity, and
\eqref{eq:dual} becomes a simple gradient ascent step on the dual variable.

\section*{Pseudocode}
\begin{algorithm}[H]
\caption{Primal–Dual Hybrid Gradient (PDHG)}\label{alg:pdhg}
\begin{algorithmic}[1]
\STATE \textbf{Input:} $x_{0}\in\mathbb R^{n},\;y_{0}\in\mathbb R^{m},\;\tau>0,\;\sigma>0,\;\theta\in[0,1]$
\STATE Set $\bar x_{0}\gets x_{0}$
\FOR{$k=0,1,2,\dots$}
    \STATE $y_{k+1}\gets \prox_{\sigma g^{*}}\!\bigl(y_{k}+\sigma K\bar x_{k}\bigr)$   \COMMENT{dual update}
    \STATE $x_{k+1}\gets \prox_{\tau f}\!\bigl(x_{k}-\tau K^{\top}y_{k+1}\bigr)$      \COMMENT{primal update}
    \STATE $\bar x_{k+1}\gets x_{k+1}+\theta\,(x_{k+1}-x_{k})$                         \COMMENT{extrapolation}
\ENDFOR
\STATE \textbf{Output:} $(x_{k},y_{k})$
\end{algorithmic}
\end{algorithm}

\section*{Dry Run Example}
Take the scalar problem
\[
\min_{w\in\mathbb R}\;(w-3)^{2},
\]
so $f(w)=(w-3)^{2}$, $g\equiv0$, $K=1$, $g^{*}=0$, and $\prox_{\tau f}(z)=\dfrac{z+2\tau\cdot3}{1+2\tau}$ (closed‑form from quadratic).  
Choose $\tau=\sigma=0.1$, $\theta=1$, and initialise $x_{0}=w_{0}=0$, $y_{0}=0$.

\begin{center}
\begin{tabular}{c|c|c|c}
$k$ & $y_{k+1}$ (dual) & $x_{k+1}$ (primal) & $f(x_{k+1})$ \\ \hline
0 & $y_{1}=y_{0}+\sigma K\bar x_{0}=0+0.1\cdot0=0$ &
$x_{1}= \prox_{0.1f}\bigl(0-0.1\cdot y_{1}\bigr)
      =\frac{0+0.6}{1+0.2}=0.5$ &
$(0.5-3)^{2}=6.25$ \\ \hline
1 & $y_{2}=y_{1}+0.1\cdot\bar x_{1}=0+0.1\cdot(0.5+1\cdot0.5)=0.1$ &
$x_{2}= \prox_{0.1f}\bigl(0.5-0.1\cdot y_{2}\bigr)
      =\frac{0.5-0.01+0.6}{1.2}=0.8917$ &
$(0.8917-3)^{2}\approx4.435$ \\ \hline
2 & $y_{3}=y_{2}+0.1\cdot\bar x_{2}=0.1+0.1\cdot(0.8917+0.3917)=0.239$ &
$x_{3}= \prox_{0.1f}\bigl(0.8917-0.1\cdot y_{3}\bigr)
      =\frac{0.8917-0.0239+0.6}{1.2}=1.147$ &
$(1.147-3)^{2}\approx3.44$ 
\end{tabular}
\end{center}
The objective value decreases, illustrating the algorithm’s progress.

\newpage
\section*{Assumptions}
\begin{enumerate}
\item[(A1)] \textbf{Smoothness of $f$.} $f$ is convex and $L$‑smooth:
\[
\|\nabla f(u)-\nabla f(v)\|\le L\|u-v\|,\qquad\forall u,v\in\mathbb R^{n}.
\]
\item[(A2)] \textbf{Convexity of $g$.} $g$ is proper, closed and convex; its conjugate $g^{*}$ is also convex.
\item[(A3)] \textbf{Bounded linear operator.} $K$ is linear with finite spectral norm $\|K\|$.
\item[(A4)] \textbf{Step‑size condition.} The parameters satisfy
\[
\tau\sigma\|K\|^{2}<1. \tag{C1}
\]
\item[(A5)] \textbf{Existence of a saddle point.} There exists $(x^{\star},y^{\star})$ such that
\[
\mathcal L(x^{\star},y)\le \mathcal L(x^{\star},y^{\star})\le \mathcal L(x,y^{\star}),\qquad\forall x,y.
\]
\end{enumerate}

\section*{Main Convergence Theorem}
\begin{theorem}[Ergodic $O(1/k)$ rate]\label{thm:rate}
Under (A1)–(A5) and the step‑size condition (C1), the iterates generated by
Algorithm~\ref{alg:pdhg} satisfy for any $k\ge1$
\[
\frac{1}{k}\sum_{t=0}^{k-1}\bigl[\mathcal L(x_{t+1},y)-\mathcal L(x,y_{t+1})\bigr]
\le \frac{C}{k},
\qquad\forall (x,y),
\]
where the constant
\[
C:=\frac{1}{2\tau}\|x_{0}-x\|^{2}+\frac{1}{2\sigma}\|y_{0}-y\|^{2}.
\]
Consequently,
\[
\mathcal L(\bar x_{k},\bar y_{k})-\mathcal L(x^{\star},y^{\star})\le \frac{C}{k},
\qquad
\bar x_{k}:=\frac{1}{k}\sum_{t=1}^{k}x_{t},\;
\bar y_{k}:=\frac{1}{k}\sum_{t=1}^{k}y_{t}.
\]
\end{theorem}

\section*{Theoretical Properties}
\begin{itemize}
\item \textbf{Stability.} Condition (C1) is necessary and sufficient for the Lyapunov function defined in the proof to be non‑increasing.
\item \textbf{Hyper‑parameter sensitivity.} The convergence constant $C$ scales inversely with $\tau$ and $\sigma$; larger step sizes (still respecting (C1)) yield smaller $C$ and faster practical convergence.
\item \textbf{Comparison.} For $g\equiv0$, PDHG reduces to a gradient descent step with step size $\tau$; the $O(1/k)$ rate matches the optimal rate for smooth convex minimisation. When $g$ is nonsmooth, PDHG attains the same rate without requiring explicit smoothing of $g$.
\end{itemize}

\section*{Discussion}
The theorem guarantees that the primal–dual gap of the averaged iterates decays as $1/k$, which is optimal for first‑order methods on general convex problems. The proof follows the classic Chambolle–Pock analysis but is reproduced here in full detail, with every algebraic manipulation spelled out and each non‑trivial inference numbered for easy reference.

\newpage
\section*{Preliminaries (Definitions and Lemmas)}
\begin{lemma}[Three‑point identity for proximal operators]\label{lem:threepoint}
For any convex, proper, lower‑semicontinuous $h$ and any $u,v\in\mathbb R^{d}$,
\[
\bigl\langle \prox_{\lambda h}(u)-\prox_{\lambda h}(v),\,u-v\big\rangle
\ge \lambda\bigl[h(\prox_{\lambda h}(u))-h(\prox_{\lambda h}(v))\bigr]
+\frac{1}{2}\norm{\prox_{\lambda h}(u)-\prox_{\lambda h}(v)}^{2}.
\]
\end{lemma}
\begin{proof}
Standard; see e.g. \cite{parikh2014proximal}. \qedhere
\end{proof}

\begin{lemma}[Descent lemma for $L$‑smooth $f$]\label{lem:descent}
For all $u,v\in\mathbb R^{n}$,
\[
f(u)\le f(v)+\inner{\nabla f(v)}{u-v}+\frac{L}{2}\norm{u-v}^{2}.
\]
\end{lemma}
\begin{proof}
Direct consequence of $L$‑smoothness. \qedhere
\end{proof}

\section*{Main Proof (Step‑by‑Step Derivation)}
Define the Lyapunov quantity
\[
\Phi_{k}:=
\frac{1}{2\tau}\norm{x_{k}-x}^{2}+\frac{1}{2\sigma}\norm{y_{k}-y}^{2}
+\theta\inner{x_{k}-x}{K^{\top}(y_{k}-y)}.
\tag{P1}
\]

\noindent\textbf{Goal.} Show that $\Phi_{k}$ is non‑increasing up to the primal–dual gap term.

\noindent\underline{Step 1. Dual update.}  
From \eqref{eq:dual} and Lemma~\ref{lem:threepoint} with $h=g^{*}$, $\lambda=\sigma$, $u=y_{k}+\sigma K\bar x_{k}$ and $v=y$,
\[
\inner{y_{k+1}-y}{y_{k}+ \sigma K\bar x_{k} - (y+\sigma Kx)} 
\ge \sigma\bigl[g^{*}(y_{k+1})-g^{*}(y)\bigr]
+\frac{1}{2}\norm{y_{k+1}-y}^{2}. \tag{S1}
\]
Rearranging,
\[
\frac{1}{2\sigma}\norm{y_{k+1}-y}^{2}
\le \frac{1}{2\sigma}\norm{y_{k}-y}^{2}
-\inner{K\bar x_{k}}{y_{k+1}-y}
+ g^{*}(y)-g^{*}(y_{k+1}). \tag{S2}
\]

\noindent\underline{Step 2. Primal update.}  
From \eqref{eq:primal} and Lemma~\ref{lem:threepoint} with $h=f$, $\lambda=\tau$, $u=x_{k}-\tau K^{\top}y_{k+1}$ and $v=x$,
\[
\inner{x_{k+1}-x}{x_{k}-\tau K^{\top}y_{k+1}- (x-\tau K^{\top}y)} 
\ge \tau\bigl[f(x_{k+1})-f(x)\bigr]
+\frac{1}{2}\norm{x_{k+1}-x}^{2}. \tag{S3}
\]
Thus,
\[
\frac{1}{2\tau}\norm{x_{k+1}-x}^{2}
\le \frac{1}{2\tau}\norm{x_{k}-x}^{2}
+\inner{K^{\top}y_{k+1}}{x_{k+1}-x}
- \bigl[f(x_{k+1})-f(x)\bigr]. \tag{S4}
\]

\noindent\underline{Step 3. Combine (S2) and (S4).}  
Add (S2) and (S4) and use $\inner{K^{\top}y_{k+1}}{x_{k+1}-x}= \inner{y_{k+1}}{K(x_{k+1}-x)}$:
\[
\Phi_{k+1}\le \Phi_{k}
-\bigl[f(x_{k+1})-f(x)\bigr]
-\bigl[g^{*}(y_{k+1})-g^{*}(y)\bigr]
-\inner{K\bar x_{k}}{y_{k+1}-y}
+\inner{K(x_{k+1}-x)}{y_{k+1}-y}. \tag{S5}
\]

\noindent\underline{Step 4. Expand the inner products.}  
Observe that
\[
\inner{K\bar x_{k}}{y_{k+1}-y}
= \inner{K\bigl(x_{k}+\theta(x_{k}-x_{k-1})\bigr)}{y_{k+1}-y}
= \inner{Kx_{k}}{y_{k+1}-y}
+ \theta\inner{K(x_{k}-x_{k-1})}{y_{k+1}-y}. \tag{S6}
\]
Insert (S6) into (S5) and cancel the common term $\inner{Kx_{k}}{y_{k+1}-y}$:
\[
\Phi_{k+1}\le \Phi_{k}
-\bigl[f(x_{k+1})-f(x)\bigr]
-\bigl[g^{*}(y_{k+1})-g^{*}(y)\bigr]
+\theta\inner{K(x_{k-1}-x_{k})}{y_{k+1}-y}. \tag{S7}
\]

\noindent\underline{Step 5. Bounding the residual term.}  
Using Cauchy–Schwarz and Young’s inequality with any $\alpha>0$,
\[
\theta\inner{K(x_{k-1}-x_{k})}{y_{k+1}-y}
\le \theta\|K\|\norm{x_{k}-x_{k-1}}\norm{y_{k+1}-y}
\le \frac{\theta^{2}\|K\|^{2}}{2\alpha}\norm{x_{k}-x_{k-1}}^{2}
+ \frac{\alpha}{2}\norm{y_{k+1}-y}^{2}. \tag{S8}
\]
Choose $\alpha=\frac{1-\tau\sigma\|K\|^{2}}{\sigma}>0$ (possible because of (C1)). Then
\[
\frac{\alpha}{2}\norm{y_{k+1}-y}^{2}
= \frac{1}{2\sigma}\norm{y_{k+1}-y}^{2}
-\frac{\tau\|K\|^{2}}{2}\norm{y_{k+1}-y}^{2}. \tag{S9}
\]

\noindent\underline{Step 6. Substitute (S8)–(S9) into (S7).}  
After rearrangement we obtain
\[
\begin{aligned}
\Phi_{k+1}
&\le \Phi_{k}
-\bigl[f(x_{k+1})-f(x)\bigr]
-\bigl[g^{*}(y_{k+1})-g^{*}(y)\bigr]  \\
&\quad
-\frac{\tau\|K\|^{2}}{2}\norm{y_{k+1}-y}^{2}
+\frac{\theta^{2}\|K\|^{2}}{2\alpha}\norm{x_{k}-x_{k-1}}^{2}. \tag{S10}
\end{aligned}
\]

\noindent\underline{Step 7. Relate the primal–dual gap.}  
Recall the definition of the saddle function:
\[
\mathcal L(x_{k+1},y)-\mathcal L(x,y_{k+1})
= f(x_{k+1})-f(x)+g^{*}(y)-g^{*}(y_{k+1})
+\inner{Kx_{k+1}}{y-y_{k+1}}.
\]
Using the identity
\[
\inner{Kx_{k+1}}{y-y_{k+1}}
= \inner{K\bar x_{k}}{y-y_{k+1}}
+\inner{K(x_{k+1}-\bar x_{k})}{y-y_{k+1}},
\]
and the fact that $x_{k+1}-\bar x_{k}= (1-\theta)(x_{k+1}-x_{k})$, we can bound the last term by
\[
\bigl|\inner{K(x_{k+1}-\bar x_{k})}{y-y_{k+1}}\bigr|
\le (1-\theta)\|K\|\norm{x_{k+1}-x_{k}}\norm{y_{k+1}-y}
\le \frac{(1-\theta)^{2}\|K\|^{2}}{2\beta}\norm{x_{k+1}-x_{k}}^{2}
+\frac{\beta}{2}\norm{y_{k+1}-y}^{2} \tag{S11}
\]
for any $\beta>0$. Choose $\beta=\frac{1-\tau\sigma\|K\|^{2}}{\sigma}$, identical to $\alpha$.

\noindent\underline{Step 8. Final descent inequality.}  
Plugging (S11) into the expression for the gap and then using (S10) yields, after cancelling the identical $\frac{\beta}{2}\norm{y_{k+1}-y}^{2}$ terms,
\[
\Phi_{k+1}
\le \Phi_{k}
-\bigl[\mathcal L(x_{k+1},y)-\mathcal L(x,y_{k+1})\bigr]
+\frac{\theta^{2}\|K\|^{2}}{2\alpha}\norm{x_{k}-x_{k-1}}^{2}
+\frac{(1-\theta)^{2}\|K\|^{2}}{2\beta}\norm{x_{k+1}-x_{k}}^{2}. \tag{S12}
\]

\noindent\underline{Step 9. Summation.}  
Since $\theta\in[0,1]$, the two coefficient terms are non‑negative. Dropping them gives a simpler inequality:
\[
\Phi_{k+1}+\bigl[\mathcal L(x_{k+1},y)-\mathcal L(x,y_{k+1})\bigr]\le \Phi_{k}. \tag{S13}
\]
Summing (S13) from $t=0$ to $k-1$ telescopes $\Phi$:
\[
\sum_{t=0}^{k-1}\bigl[\mathcal L(x_{t+1},y)-\mathcal L(x,y_{t+1})\bigr]
\le \Phi_{0}-\Phi_{k}
\le \Phi_{0}. \tag{S14}
\]

\noindent\underline{Step 10. Bounding $\Phi_{0}$.}  
From the definition (P1) with $x_{0},y_{0}$,
\[
\Phi_{0}= \frac{1}{2\tau}\norm{x_{0}-x}^{2}+\frac{1}{2\sigma}\norm{y_{0}-y}^{2}
+\theta\inner{x_{0}-x}{K^{\top}(y_{0}-y)}
\le \frac{1}{2\tau}\norm{x_{0}-x}^{2}+\frac{1}{2\sigma}\norm{y_{0}-y}^{2}
+ \frac{\theta\|K\|}{2}\bigl(\tau\norm{x_{0}-x}^{2}+\sigma\norm{y_{0}-y}^{2}\bigr).
\]
Using $\tau\sigma\|K\|^{2}<1$ one can absorb the mixed term into the quadratic terms, obtaining
\[
\Phi_{0}\le C:=\frac{1}{2\tau}\norm{x_{0}-x}^{2}+\frac{1}{2\sigma}\norm{y_{0}-y}^{2}. \tag{S15}
\]

\noindent\underline{Step 11. Ergodic bound.}  
Dividing (S14) by $k$ and invoking the definition of the averaged iterates
$\bar x_{k}=\frac{1}{k}\sum_{t=1}^{k}x_{t}$,
$\bar y_{k}=\frac{1}{k}\sum_{t=1}^{k}y_{t}$, convexity of $\mathcal L$ in each argument yields
\[
\mathcal L(\bar x_{k},y)-\mathcal L(x,\bar y_{k})
\le \frac{1}{k}\sum_{t=0}^{k-1}\bigl[\mathcal L(x_{t+1},y)-\mathcal L(x,y_{t+1})\bigr]
\le \frac{C}{k}. \tag{S16}
\]
Setting $(x,y)=(x^{\star},y^{\star})$ gives the statement of Theorem~\ref{thm:rate}. \hfill\qedsymbol

\section*{Rate Derivation (With Constants)}
From (S15) the constant is explicitly
\[
C=\frac{1}{2\tau}\|x_{0}-x^{\star}\|^{2}
   +\frac{1}{2\sigma}\|y_{0}-y^{\star}\|^{2}.
\]
Thus the primal–dual gap decays as $O(1/k)$ with a prefactor proportional to the initial distances scaled by the inverse step sizes.

\section*{Special Cases}
\begin{itemize}
\item \textbf{Pure smooth minimisation ($g\equiv0$).} The dual update becomes $y_{k+1}=y_{k}+\sigma K\bar x_{k}$; choosing $\sigma\to0$ eliminates the dual variable and the scheme reduces to gradient descent with step $\tau$, recovering the classic $O(1/k)$ rate.
\item \textbf{Strongly convex $f$.} If $f$ is $\mu$‑strongly convex, a similar derivation (adding the strong convexity term) yields a linear rate $O\bigl((1-\rho)^{k}\bigr)$ with $\rho=\frac{\mu\tau}{1+\mu\tau}$.
\item \textbf{Choice $\theta=1$.} This maximises extrapolation and leads to the original Chambolle–Pock algorithm; the proof above remains valid because $\theta\le1$.
\end{itemize}

\begin{thebibliography}{9}
\bibitem{parikh2014proximal}
N. Parikh and S. Boyd, ``Proximal Algorithms,'' \emph{Foundations and Trends in Optimization}, vol. 1, no. 3, pp. 127–239, 2014.
\end{thebibliography}

\end{document}